\documentclass[aspectratio=169,11pt]{beamer}

\usepackage[utf8]{inputenc}
\usepackage[T1]{fontenc}
\usepackage[brazil]{babel}
\usepackage{graphicx}
\usepackage{xcolor}
\usepackage{booktabs}
\usepackage{tikz}

\usetikzlibrary{positioning}

\definecolor{Ink}{HTML}{17202A}
\definecolor{Panel}{HTML}{F6F8FA}
\definecolor{RoadBlue}{HTML}{28A0FF}
\definecolor{RoadOrange}{HTML}{DC9128}
\definecolor{SignalRed}{HTML}{E84A5F}
\definecolor{SoftGreen}{HTML}{2F7D5C}

\setbeamercolor{normal text}{fg=Ink,bg=white}
\setbeamercolor{frametitle}{fg=Ink,bg=white}
\setbeamercolor{title}{fg=Ink}
\setbeamercolor{structure}{fg=RoadBlue}
\setbeamertemplate{navigation symbols}{}
\setbeamertemplate{footline}{%
    \hfill\color{Ink!55}\scriptsize\insertframenumber/\inserttotalframenumber\hspace{0.45cm}\vspace{0.25cm}
}
\setbeamertemplate{frametitle}{%
    \vspace{0.28cm}
    \begin{beamercolorbox}[wd=\paperwidth,leftskip=0.55cm,rightskip=0.55cm]{frametitle}
        \usebeamerfont{frametitle}\insertframetitle
    \end{beamercolorbox}
    \vspace{-0.15cm}
    \hspace{0.55cm}\textcolor{RoadBlue}{\rule{1.15cm}{1.5pt}}\textcolor{RoadOrange}{\rule{0.55cm}{1.5pt}}
}

\newcommand{\resultimage}[2]{%
    \IfFileExists{data/results/#1}{%
        \includegraphics[#2]{data/results/#1}%
    }{%
        \IfFileExists{../data/results/#1}{%
            \includegraphics[#2]{../data/results/#1}%
        }{%
            \fbox{\parbox[c][4.4cm][c]{0.86\textwidth}{\centering
            Execute \texttt{make run} para gerar \texttt{#1}.}}%
        }%
    }%
}

\newcommand{\pill}[2]{%
    \tikz[baseline=(n.base)]\node[
        fill=#1!12,
        draw=#1!55,
        rounded corners=2pt,
        inner xsep=7pt,
        inner ysep=3pt,
        text=#1!55!Ink
    ] (n) {\scriptsize #2};
}

\title{Detecção e Reconstrução de Malha Viária}
\subtitle{Reconhecimento de Padrões em Imagens de Satélite}
\author{Beatriz Vocurca Frade \and Johnatan Augusto Moreira do Carmo \and Marina Alves Resende}
\date{\today}

\begin{document}

\begin{frame}[plain]
    \begin{tikzpicture}[remember picture,overlay]
        \node[anchor=east] at ([xshift=0.05cm]current page.east) {
            \resultimage{demo_sintetico_sobreposicao.png}{height=\paperheight}
        };
        \fill[white,opacity=0.94] (current page.south west) rectangle ([xshift=0.57\paperwidth]current page.north west);
        \fill[RoadBlue] ([xshift=0.35cm,yshift=-0.35cm]current page.north west) rectangle ++(0.14,0.14);
        \fill[RoadOrange] ([xshift=0.54cm,yshift=-0.35cm]current page.north west) rectangle ++(0.14,0.14);
    \end{tikzpicture}
    \vspace{1.1cm}
    {\Huge\bfseries Detecção e\\Reconstrução de\\Malha Viária\par}
    \vspace{0.35cm}
    {\large Imagens de satélite + visão computacional clássica\par}
    \vspace{0.75cm}
    \pill{RoadBlue}{segmentação}
    \pill{RoadOrange}{morfologia}
    \pill{SignalRed}{grafos}
    \vfill
    {\small Beatriz Vocurca Frade \quad Johnatan Augusto Moreira do Carmo \quad Marina Alves Resende}
\end{frame}

\begin{frame}{Problema}
    \begin{columns}[T,totalwidth=\textwidth]
        \begin{column}{0.48\textwidth}
            \Large Identificar vias e reconstruir sua conectividade a partir de uma imagem.
            \vspace{0.45cm}

            \begin{itemize}
                \item vias asfaltadas e de terra;
                \item cruzamentos e extremidades;
                \item sobreposição visual;
                \item grafo exportável.
            \end{itemize}
        \end{column}
        \begin{column}{0.48\textwidth}
            \begin{tikzpicture}
                \node[inner sep=0pt] (img) {\resultimage{demo_sintetico_sobreposicao.png}{width=\linewidth}};
                \draw[RoadBlue,line width=2pt] (img.south west) -- ++(1.35,0);
                \draw[RoadOrange,line width=2pt] ([xshift=1.45cm]img.south west) -- ++(1.35,0);
            \end{tikzpicture}
        \end{column}
    \end{columns}
\end{frame}

\begin{frame}{Desafios}
    \begin{columns}[T,totalwidth=\textwidth]
        \begin{column}{0.31\textwidth}
            \pill{RoadBlue}{aparência variável}
            \vspace{0.25cm}

            Mesma classe, cores e texturas diferentes conforme sensor, iluminação, escala e região.
        \end{column}
        \begin{column}{0.31\textwidth}
            \pill{RoadOrange}{falsos positivos}
            \vspace{0.25cm}

            Telhados, solo exposto, estacionamentos e sombras podem parecer vias.
        \end{column}
        \begin{column}{0.31\textwidth}
            \pill{SignalRed}{topologia frágil}
            \vspace{0.25cm}

            Pequenas quebras na máscara podem quebrar arestas ou criar cruzamentos falsos.
        \end{column}
    \end{columns}

    \vspace{0.75cm}
    \begin{center}
        \begin{tikzpicture}[node distance=0.25cm]
            \node[draw=Ink!15,fill=Panel,rounded corners=3pt,minimum width=2.15cm,minimum height=0.9cm] (a) {imagem};
            \node[draw=Ink!15,fill=Panel,rounded corners=3pt,minimum width=2.15cm,minimum height=0.9cm,right=of a] (b) {máscara};
            \node[draw=Ink!15,fill=Panel,rounded corners=3pt,minimum width=2.15cm,minimum height=0.9cm,right=of b] (c) {esqueleto};
            \node[draw=Ink!15,fill=Panel,rounded corners=3pt,minimum width=2.15cm,minimum height=0.9cm,right=of c] (d) {grafo};
            \draw[RoadBlue,line width=1.5pt,->] (a) -- (b);
            \draw[RoadBlue,line width=1.5pt,->] (b) -- (c);
            \draw[RoadBlue,line width=1.5pt,->] (c) -- (d);
        \end{tikzpicture}
    \end{center}
\end{frame}

\begin{frame}{Pipeline}
    \begin{center}
        \begin{tikzpicture}[node distance=0.23cm]
            \node[draw=RoadBlue!30,fill=RoadBlue!8,rounded corners=3pt,minimum width=2.35cm,minimum height=0.95cm] (a) {pré-processamento};
            \node[draw=RoadBlue!30,fill=RoadBlue!8,rounded corners=3pt,minimum width=2.35cm,minimum height=0.95cm,right=of a] (b) {descritores};
            \node[draw=RoadOrange!45,fill=RoadOrange!10,rounded corners=3pt,minimum width=2.35cm,minimum height=0.95cm,right=of b] (c) {segmentação};
            \node[draw=RoadOrange!45,fill=RoadOrange!10,rounded corners=3pt,minimum width=2.35cm,minimum height=0.95cm,right=of c] (d) {morfologia};
            \node[draw=SignalRed!45,fill=SignalRed!8,rounded corners=3pt,minimum width=2.35cm,minimum height=0.95cm,right=of d] (e) {grafo};
            \draw[->,line width=1.4pt,Ink!45] (a) -- (b);
            \draw[->,line width=1.4pt,Ink!45] (b) -- (c);
            \draw[->,line width=1.4pt,Ink!45] (c) -- (d);
            \draw[->,line width=1.4pt,Ink!45] (d) -- (e);
        \end{tikzpicture}
    \end{center}

    \vspace{0.55cm}
    \begin{columns}[T,totalwidth=\textwidth]
        \begin{column}{0.48\textwidth}
            \begin{itemize}
                \item filtro bilateral e CLAHE;
                \item RGB, HSV, Lab e tons de cinza;
                \item textura local e suporte linear.
            \end{itemize}
        \end{column}
        \begin{column}{0.48\textwidth}
            \begin{itemize}
                \item regras explicáveis;
                \item esqueletonização;
                \item nós, arestas, JSON e GraphML.
            \end{itemize}
        \end{column}
    \end{columns}
\end{frame}

\begin{frame}{Segmentação}
    \begin{columns}[T,totalwidth=\textwidth]
        \begin{column}{0.48\textwidth}
            \resultimage{demo_sintetico_segmentada.png}{width=\linewidth}
        \end{column}
        \begin{column}{0.48\textwidth}
            \Large Regras combinadas, não um único limiar.
            \vspace{0.35cm}

            \begin{itemize}
                \item \textcolor{RoadBlue}{azul}: asfalto;
                \item \textcolor{RoadOrange}{laranja}: terra;
                \item remoção de vegetação e sombras;
                \item solo de terra exige suporte linear.
            \end{itemize}
        \end{column}
    \end{columns}
\end{frame}

\begin{frame}{Sobreposição}
    \begin{columns}[T,totalwidth=\textwidth]
        \begin{column}{0.58\textwidth}
            \resultimage{demo_sintetico_sobreposicao.png}{height=0.72\textheight}
        \end{column}
        \begin{column}{0.38\textwidth}
            \Large A sobreposição preserva o contexto original.
            \vspace{0.35cm}

            \begin{itemize}
                \item facilita inspeção visual;
                \item evidencia falsos positivos;
                \item mostra continuidade das vias.
            \end{itemize}
        \end{column}
    \end{columns}
\end{frame}

\begin{frame}{Esqueleto}
    \begin{columns}[T,totalwidth=\textwidth]
        \begin{column}{0.50\textwidth}
            \resultimage{demo_sintetico_esqueleto.png}{width=\linewidth}
        \end{column}
        \begin{column}{0.46\textwidth}
            \Large A máscara vira uma estrutura topológica.
            \vspace{0.35cm}

            \begin{itemize}
                \item grau 1: extremidade;
                \item grau 2: trecho intermediário;
                \item grau 3 ou mais: interseção;
                \item caminhos entre nós viram arestas.
            \end{itemize}
        \end{column}
    \end{columns}
\end{frame}

\begin{frame}{Grafo}
    \begin{columns}[T,totalwidth=\textwidth]
        \begin{column}{0.58\textwidth}
            \resultimage{demo_sintetico_grafo.png}{height=0.72\textheight}
        \end{column}
        \begin{column}{0.38\textwidth}
            \begin{table}
                \small
                \begin{tabular}{lr}
                    \toprule
                    Métrica & Valor \\
                    \midrule
                    Pixels de via & 44.037 \\
                    Componentes & 31 \\
                    Nós & 18 \\
                    Arestas & 19 \\
                    Asfalto & 24.469 \\
                    Terra & 19.568 \\
                    \bottomrule
                \end{tabular}
            \end{table}
            \vspace{0.2cm}
            \pill{SignalRed}{JSON + GraphML}
        \end{column}
    \end{columns}
\end{frame}

\begin{frame}{Limitações}
    \begin{columns}[T,totalwidth=\textwidth]
        \begin{column}{0.48\textwidth}
            \Large O método é explicável, mas não universal.
            \vspace{0.3cm}

            \begin{itemize}
                \item sem base anotada, avaliação é qualitativa;
                \item objetos parecidos com vias ainda geram ruído;
                \item o grafo depende da continuidade da máscara.
            \end{itemize}
        \end{column}
        \begin{column}{0.48\textwidth}
            \Large Próximos passos.
            \vspace{0.3cm}

            \begin{itemize}
                \item medir IoU, precisão, revocação e F1;
                \item testar U-Net, DeepLab ou SegFormer;
                \item comparar com OpenStreetMap;
                \item refinar conectividade topológica.
            \end{itemize}
        \end{column}
    \end{columns}
\end{frame}

\begin{frame}{Conclusão}
    \begin{center}
        {\Huge\bfseries Visão computacional clássica\\com saída topológica}
    \end{center}

    \vspace{0.45cm}
    \begin{columns}[T,totalwidth=\textwidth]
        \begin{column}{0.32\textwidth}
            \pill{RoadBlue}{segmenta}

            Identifica pixels candidatos a vias.
        \end{column}
        \begin{column}{0.32\textwidth}
            \pill{RoadOrange}{interpreta}

            Separa asfalto e terra para visualização.
        \end{column}
        \begin{column}{0.32\textwidth}
            \pill{SignalRed}{reconstrói}

            Exporta nós, arestas e conectividade.
        \end{column}
    \end{columns}

    \vfill
    \begin{center}
        \texttt{make run} \quad \texttt{make report} \quad \texttt{make slides}
    \end{center}
\end{frame}

\end{document}
